\section{The contact line and the Picard alternatives}
\label{red:section}

Let $X$ be a smooth connected complex projective contact Fano manifold
of dimension $d=2n+1$, where $n\geq1$, with contact sequence
\begin{equation}\label{red:contact}
0\longrightarrow F\longrightarrow T_X
\xrightarrow{\theta}L\longrightarrow0.
\end{equation}
In a local frame of $L$ we write $\alpha$ for $\theta$. With the Levi
convention in Section~\ref{sec:introduction}, local sections $v,w$ of
$F$ satisfy $d\alpha(v,w)=-\alpha([v,w])$.
The top exterior power of $d\alpha|_F$ is nowhere zero. Taking determinants gives
\begin{equation}\label{red:determinant}
\det F\simeq L^{\otimes n},\qquad
-K_X\simeq L^{\otimes(n+1)}.
\end{equation}
In particular $L$ is ample.

\begin{lemma}\label{red:simply-connected}
The manifold $X$ is simply connected, and its Picard group is a free
abelian group of finite rank, canonically isomorphic to $H^2(X,\ZZ)$.
\end{lemma}

\begin{proof}
We recall an argument that applies to every smooth complex Fano
manifold. Yau's prescribed-Ricci-form Theorem~\cite[Theorem~1]{Yau1977},
proved in~\cite{Yau1978}, says that
a real closed $(1,1)$-form representing $2\pi c_1(X)$ is the Ricci form
of a K\"ahler metric in any prescribed K\"ahler class. Ampleness of
$-K_X$ supplies a positive representative. For the resulting metric,
compactness gives a uniform positive lower bound for the Ricci tensor.
The Bonnet--Myers Theorem~\cite{Myers1941} makes the universal cover
$q:\widetilde X\to X$ compact, so $q$ has finite degree $r$.

The cover is holomorphic and unramified, and the lifted metric
is K\"ahler with positive Ricci tensor. In particular
$-K_{\widetilde X}=q^*(-K_X)$ is positive. Kodaira's projectivity
Theorem~\cite[Section~3, p.~316]{Kodaira1954} therefore makes the
compact manifold $\widetilde X$ projective. The graph of $q$ is a
closed analytic subvariety of the projective product
$\widetilde X\times X$, so GAGA makes $q$ algebraic
\cite[Proposition~15]{Serre1956}. It is proper with finite fibers, hence finite,
and its local analytic inverses show that it is \'etale.
Thus both manifolds and their finite covering are algebraic before
we apply projective vanishing and Riemann--Roch.

Kodaira vanishing~\cite{Kodaira1953}, applied to $\OO=K\otimes K^{-1}$, gives
$H^i(\widetilde X,\OO_{\widetilde X})=H^i(X,\OO_X)=0$ for $i>0$.
Both holomorphic Euler characteristics are one. Since
$T_{\widetilde X}=q^*T_X$, Hirzebruch--Riemann--Roch~\cite{Hirzebruch1954} gives
\[
1=\chi(\widetilde X,\OO_{\widetilde X})
 =\int_{\widetilde X}q^*\operatorname{td}(T_X)
 =r\,\chi(X,\OO_X)=r.
\]
Thus $X$ is simply connected. This argument prescribes a Ricci form;
it does not require the metric to be K\"ahler--Einstein.

The exponential sequence and the vanishings of $H^1(X,\OO_X)$ and
$H^2(X,\OO_X)$ identify the analytic Picard group with $H^2(X,\ZZ)$;
GAGA~\cite{Serre1956} identifies it with the algebraic Picard group. The cohomology
group is finitely generated. The universal coefficient sequence
has torsion term $\operatorname{Ext}^1(H_1(X,\ZZ),\ZZ)=0$, since
$H_1(X,\ZZ)=0$, and identifies $H^2(X,\ZZ)$ with a subgroup of a
finitely generated free abelian group. It is therefore free.
\end{proof}

\begin{proposition}\label{red:alternatives}
Either $\operatorname{Pic}(X)=\ZZ[L]$, or $(X,F)$ is contact-isomorphic
to one of the following standard contact manifolds:
\begin{enumerate}
\item the incidence variety
\[
I=\{(\ell,H)\in\PP^{n+1}\times(\PP^{n+1})^*: \ell\subset H\},
\qquad L=\OO_I(1,1);
\]
\item projective space $\PP^{2n+1}$ with the contact structure induced
by a nondegenerate alternating form on $\CC^{2n+2}$, and
$L=\OO_{\PP^{2n+1}}(2)$.
\end{enumerate}
In both listed cases a connected algebraic group acts transitively
and preserves the given distribution $F$.
\end{proposition}

\begin{proof}
Suppose first that $\operatorname{rank}\operatorname{Pic}(X)>1$.
Lemma~\ref{red:simply-connected} gives $b_2(X)>1$, and $K_X$ is not
nef. Theorem~1.1 of Kebekus--Peternell--Sommese--Wi\'sniewski
\cite{KPSW2000} states that a projective contact manifold with
non-nef canonical bundle is either Fano with $b_2=1$, or a
projectivized tangent bundle over a smooth projective manifold.
Their convention is the projectivization of one-dimensional
quotients. In our convention of lines, the second alternative is
\[
X\simeq\PP(T^*Z),\qquad \dim Z=n+1.
\]
For a rank-$r$ bundle $E$ in the lines convention,
\[
K_{\PP(E)}=\OO_{\PP(E)}(-r)
 \otimes\pi^*(K_Z\otimes(\det E)^{-1}).
\]
Taking $E=T^*Z$ cancels the factor pulled back from $Z$, so
$-K_X=\OO_{\PP(T^*Z)}(n+1)$. Consequently
$\OO_{\PP(T^*Z)}(1)$ is ample. This is precisely ampleness of $T_Z$,
as defined by its tautological quotient line bundle. Mori's
Theorem~\cite{Mori1979}, that a smooth complex projective manifold
with ample tangent bundle is projective space, now gives
$Z\simeq\PP^{n+1}$. Thus $X\simeq I$.

We verify the contact structure, rather than only the variety.
Put $m=n+1\geq2$ and
$A=\PP^m\times(\PP^m)^*$. The incidence variety is a smooth
divisor of type $(1,1)$, so adjunction gives
$-K_I=\OO_I(m,m)$. Equation~\eqref{red:determinant} and
torsion-freeness of $\operatorname{Pic}(I)$ give $L=\OO_I(1,1)$.
Set $E=\Omega_A^1(1,1)|_I$. The restriction sequence is
\[
0\longrightarrow\Omega_A^1\longrightarrow\Omega_A^1(1,1)
 \longrightarrow E\longrightarrow0.
\]
The Euler sequence on $\PP^m$ gives
$H^i(\PP^m,\Omega_{\PP^m}^1(1))=0$ for all $i$ and
$h^1(\PP^m,\Omega_{\PP^m}^1)=1$.
The decomposition of $\Omega_A^1$ and the K\"unneth formula imply
\[
H^0(A,\Omega_A^1(1,1))=H^1(A,\Omega_A^1(1,1))=0,
\qquad h^1(A,\Omega_A^1)=2.
\]
Hence $h^0(I,E)=2$. The twisted conormal sequence
\[
0\longrightarrow\OO_I\longrightarrow E
 \longrightarrow\Omega_I^1(1,1)\longrightarrow0
\]
and $H^1(I,\OO_I)=0$ then give
$h^0(I,\Omega_I^1(1,1))=1$.
The tautological contact form on $\PP(T^*\PP^m)$ is a nonzero
section of this space. Indeed, at $(\ell,H)$ with
$\ell\subset H\subset\CC^{m+1}$, the cotangent line is
$\ell\otimes(\CC^{m+1}/H)^*$. Thus the tautological subbundle
is $\OO_I(-1,-1)$ and its dual contact line is $\OO_I(1,1)$.
The given contact form is a
nonzero scalar multiple of the tautological form, and the kernels
agree. The simultaneous action of $\operatorname{PGL}_{m+1}$ on
points and hyperplanes is transitive on $I$ and preserves this kernel.

Suppose next that $\operatorname{Pic}(X)$ has rank one. Choose its
ample generator $A_1$ and write $L=A_1^{\otimes k}$ with $k>0$.
If $k\geq2$, the Fano index is $(n+1)k\geq2n+2=d+1$.
The Kobayashi--Ochiai index Theorem~\cite{KO1973} bounds the index of
a smooth Fano $d$-fold by $d+1$, with equality only for
$\PP^d$. Thus $k=2$, $X\simeq\PP^d$, and $L=\OO(2)$.
For $X=\PP(V)$ the Euler sequence identifies
\[
H^0(X,\Omega_X^1(2))\simeq\bigwedge\nolimits^2 V^*.
\]
The twisted one-form corresponding to $\sigma\in\bigwedge^2V^*$
vanishes at $[v]$ if $v$ belongs to the radical of $\sigma$.
Our contact quotient is everywhere surjective, so its alternating
form has zero radical. All nondegenerate alternating forms on $V$
are equivalent under $\operatorname{GL}(V)$. This identifies the
given contact distribution with the standard one, preserved by
the transitive connected symplectic group.

The only remaining alternative is $k=1$.
\end{proof}

The two listed manifolds are the type $A_{n+1}$ and $C_{n+1}$ adjoint
contact varieties. In the projective-space case the adjoint
polarization is $\OO(2)$, so its adjoint embedding is the quadratic
Veronese embedding. Their contact transitivity also allows them to
be included in the uniform homogeneous identification at the end
of the proof. In the intervening sections we assume
\begin{equation}\label{red:primitive}
\operatorname{Pic}(X)=\ZZ[L].
\end{equation}
